\documentclass[10pt,a4paper]{amsart}
\usepackage[margin=19mm]{geometry}
\usepackage{amsmath,amssymb,mathtools}
\usepackage[scr=rsfs,cal=euler]{mathalfa}
\usepackage{xfrac}
\usepackage[unicode,hidelinks,bookmarks=false]{hyperref}
\newcommand{\ie}{i.e.\ }
\newcommand{\cf}{cf.\ }
\newcommand{\resp}{respectively\ }
\newcommand{\Subsection}{\subsection}
\providecommand{\nobreakdash}{\nobreak\hbox{-}}
\setlength{\emergencystretch}{2em}
\multlinegap0pt
\usepackage{bm}
\usepackage{bbm}
\usepackage{dsfont}
\usepackage{xspace}
\usepackage{cancel}
\usepackage{mathtools}
\usepackage{amsthm}
\makeatletter
\@ifundefined{theorem}{\newtheorem{theorem}{Theorem}}{}
\makeatother
\newcommand{\WF}{\mathrm{WF}}                         % Wavefront set
\newcommand{\dd}{\mathrm{d}}  % without the space
\newcommand{\vx}{{\mathbf{x}}}
\title[Microlocal analysis of non-linear operators arising in Compt]{Microlocal analysis of non-linear operators arising in Compton CT}
\author{James W. Webber, Sean Holman}
\date{Source: arXiv:2507.19791v3 (CC BY 4.0)}
\begin{document}
\maketitle

\noindent\emph{Source and licence.} Microlocal analysis of non-linear operators arising in Compton CT. Version arXiv:2507.19791v3 (CC BY 4.0), distributed under the Creative Commons Attribution 4.0 International licence, as recorded in the arXiv licence metadata for this exact version and shown on the abstract page https://arxiv.org/abs/2507.19791v3. Full rights notice: licenses/arxiv-2507.19791-CC-BY-4.0.txt.\par\smallskip

\noindent\textit{Editorial scope.} Counted unit: the Theorem whose source label is \texttt{main\_2} in the article (source0.tex lines 1396--1399, source label \texttt{main\_2}), delivered together with the article's own complete original proof of it (source0.tex lines 1400--1440, one \texttt{proof} environment of 41 lines). No mathematics is written, repaired or re-derived: statement and proof are the article's own text line for line. Required context retained uncounted: trace\_thm (lines 623--630); thm\_V\_smth (lines 1017--1020); main\_thm\_1 (lines 1320--1325); main\_1 (lines 1388--1391). Every definition, display and label of those lines is retained unchanged. Omitted from this package and not redistributed: 1--622, 631--1016, 1021--1319, 1326--1387, 1392--1395, 1441--1797. Nothing inside a retained excerpt is omitted or altered: every retained body is delivered exactly as the article's lines are written, so no omission receipt of any kind accompanies this package. Citations: the retained excerpts carry no citation command at all, so no bibliography entry of the article is transcribed and no reference is redistributed. Reference adaptation: none was needed and none was made; every reference inside the delivered unit resolves inside it, so no unresolved reference or citation is produced. Printed slips kept literal: no printed word, symbol, hypothesis or number of the counted statement or of its proof was altered. Layout and class: the article's own class is \texttt{amsart} with packages including lipsum, amssymb, amsthm, amsaddr, amsmath, tikz, graphicx, subcaption, enumitem, float, hyperref, placeins, pgfplots, fullpage; the wrapper is the lane's pinned \texttt{amsart} editorial wrapper at 10pt on A4, which restates the theorem-like environments the retained text needs and adds the article's own macro definitions that the retained text needs (\texttt{\textbackslash WF}, \texttt{\textbackslash dd}, \texttt{\textbackslash vx}), with extra wrapper packages (bm, bbm, dsfont, xspace, cancel, mathtools). The delivered numbering is the unit's own. Assets: no retained span contains a figure environment, a graphics-inclusion command, a PDF-page inclusion, a table or a TikZ picture; no image file is copied into this package, the package declares no asset dependency and no image file is redistributed with it. Licence: version arXiv:2507.19791v3 of this article is distributed under the Creative Commons Attribution 4.0 International licence, as recorded in the arXiv licence metadata for this exact version and shown on the abstract page https://arxiv.org/abs/2507.19791v3. A full rights notice is delivered at licenses/arxiv-2507.19791-CC-BY-4.0.txt. Characters: every retained excerpt is pure ASCII, so the delivered engine, XeLaTeX with its default Latin Modern text font, reports no missing character.
\begin{theorem}
\label{trace_thm}
Let $Tf(\vx') = f(\vx',0)$ be a trace operator. Then, for $\alpha > 1/2$,
\begin{equation}
T: H^{\alpha,2}(\mathbb{R}^n) \to H^{\alpha-1/2,2}(\mathbb{R}^{n-1})
\end{equation}
is a bounded map.
\end{theorem}

\begin{theorem}
\label{thm_V_smth}
Let $\phi \in [-\pi,\pi]$ be fixed, $\varphi \in H^{\lambda,\infty}(\mathbb{R}^2)$, and $f \in H^{\alpha,2}(\mathbb{R}^2)$. Then, $\widetilde{\mathcal{V}_L f }(\cdot,\phi) \in H^{\alpha + \lambda,2}(\mathbb{R}^2)$.
\end{theorem}

\begin{theorem}\label{main_thm_1}
Let $f = u\chi_{\Omega}$, where $\Omega \subset B$ is a compact domain with smooth boundary, and $u \in C^\infty(\mathbb{R}^2)$. Let $\varphi \in H^{\lambda,\infty}(\mathbb{R}^2)$ for $\lambda > 1/2$, and let $g = e^{ -\widetilde{\mathcal{V}_L}f(\cdot,\phi) }$ where $\phi \in [-\pi, \pi]$ is fixed. If $(\vx_0,\xi) \notin WF(f)$ then $fg \in H^{\beta,2}$ microlocally near $(\vx_0,\xi)$ for any $\beta < \lambda$. If we also assume $u>0$, then
\begin{equation}
\WF(fg)\backslash \WF^{1/2}(fg) = \WF(f).
\end{equation}
\end{theorem}

\begin{theorem}
\label{main_1}
Let $f = u\chi_\Omega$ be as in Theorem \ref{main_thm_1}, and let $(s_0,\theta_0)$ be such that $L = L(s_0,\theta_0)$ is not tangent to $\partial \Omega$. Let $\varphi \in H^{\lambda,\infty}(\mathbb{R}^2)$ for $\lambda > 1/2$. Then $\mathcal{R}f(\cdot,\theta_0)$ is in $H^{\beta,2}$ near $s_0$ for any $\beta < \lambda$.
\end{theorem}

\begin{theorem}
\label{main_2}
Let $f = u\chi_\Omega$ be as in Theorem \ref{main_thm_1} with $u>0$ and let $(s_0,\theta_0)$ be such that $L = L(s_0,\theta_0)$ is tangent to $\partial \Omega$ and the set of tangent points is either a line segment or a single point at which the curvature of $\partial \Omega$ is not zero. Also suppose that $\varphi \in H^{\lambda,\infty}(\mathbb{R}^2)$ with $\lambda > 1$. Then, $\mathcal{R}f(\cdot,\theta_0)$ is not in $H^{1,2}$ near $s_0$.
\end{theorem}

\begin{proof}
%\seanc{Working on this Theorem and proof. We will need to change the statement because it is not true in general.}

Let $T \subset \mathbb{R}^2$ be the set of points where $L(s_0,\theta_0)$ is tangent to $\partial \Omega$. By hypothesis, $T$ is either a single point or a line segment and in either case let $\psi \in C_c^\infty(\mathbb{R}^2)$ be a cut-off function which is equal to $1$ in a neighbourhood of $T$. Then we have
\begin{equation}\label{eq:Rdecomp}
\mathcal{R}f(s,\theta_0) = \lambda \int_{L(s,\theta_0)}w \psi f e^{-\widetilde{\mathcal{V}_{L}}f(\vx,\theta_0)} \dd l + \lambda \int_{L(s,\theta_0)}w (1-\psi) f e^{-\widetilde{\mathcal{V}_{L}}f(\vx,\theta_0)} \dd l
\end{equation}
A slight variation of the proofs of Theorems \ref{main_thm_1} and \ref{main_1} shows that the second term on the right side of \eqref{eq:Rdecomp} will be in $H^{\beta,2}$ near $s_0$ for $\beta < \lambda$. Since $\lambda > 1$, it is therefore sufficient to show that the first term on the right side of \eqref{eq:Rdecomp} is not in $H^{1,2}$ near $s_0$. For notational convenience, let us write
\[
R(s) = \lambda \int_{L(s,\theta_0)}w \psi f e^{-\widetilde{\mathcal{V}_{L}}f(\vx,\theta_0)} \dd l
\]
for this term.

Before continuing, note that since $\psi f \in H^{\alpha,2}(\mathbb{R}^2)$ for any $\alpha< 1/2$ and since $\lambda > 1$, by Theorem \ref{thm_V_smth}, $\widetilde{V_L} f(\cdot,\theta_0) \in H^{3/2+,2}(\mathbb{R}^2)$ where the plus sign indicates membership in the intersection of spaces with order higher than $3/2$. Furthermore, all derivatives of $\widetilde{V_L} f(\cdot,\theta_0)$ will be in $H^{1/2+,2}(\mathbb{R}^2)$. Sobolev embedding also implies that $\widetilde{V_L} f(\cdot,\theta_0)$ is continuous.

Now, let us consider the case when $T$ is a line segment. Without loss of generality, we suppose that $\Theta_0 = \Theta(\theta_0) = (\cos(\theta),\sin(\theta))^T$ points into $\Omega$ on $T$. By continuity of the integrand,
\[
\lim_{s\rightarrow s_0^+} R(s) - \lim_{s\rightarrow s_0^-} R(s) = \int_T u e^{-\widetilde{\mathcal{V}_{L}}f(\vx,\theta_0)} \dd l > 0.
\]
Therefore $Rf(\cdot,\theta_0)$ is not continuous at $s_0$ and so cannot be in $H^{1,2}$ near $s_0$ by Sobolev embedding.

Now consider the case when $T$ is a single point $\vx_0 = s_0 \Theta_0 + t_0 (\Theta_0)_\perp$ where the curvature of $\partial \Omega$ does not vanish. Note that
\[
\lambda \int_{L(s,\theta_0)} w \psi u e^{-\widetilde{\mathcal{V}_{L}}f(\vx,\theta_0)} \dd l = R(s) + \lambda \int_{L(s,\theta_0)} w \psi u \chi_{\Omega^c} e^{-\widetilde{\mathcal{V}_{L}}f(\vx,\theta_0)} \dd l.
\]
Since $w \psi u \in C_c^\infty(\mathbb{R}^2)$ and $e^{-\widetilde{\mathcal{V}_{L}}f(\vx,\theta_0)} \in H^{3/2+,2}(\mathbb{R}^2)$, the term on the left side of this equation is in $H^{1,2}(\mathbb{R})$. Therefore, to show $R$ is not in $H^{1,2}(\mathbb{R})$ it is sufficient to show that either term on the right side is not in $H^{1,2}(\mathbb{R})$. Thus, by possibly replacing $\Omega$ by $\Omega^c$, we can assume without loss of generality that, at $\vx_0$, $\Theta_0$ points into $\Omega$ and when $\partial \Omega$ is oriented by $\Theta_0$ its curvature is positive.

Next, taking the support of $\psi$ to be sufficiently small, for $s<s_0$ we have $R(s) = 0$ while for $s>s_0$, since the curvature is positive, there will be exactly two values $t_\pm(s)$ such that $\vx_0 + s \Theta_0 + t_\pm(s) (\Theta_0)_\perp \in \partial \Omega$ within the support of $\psi$. We suppose $\pm t_\pm(s)>0$. For $s>s_0$ in the support of $\psi$, we then have
\[
R(s) = \lambda \int_{t_-(s)}^{t_+(s)} wu(\vx_0 + s  \Theta_0 + t (\Theta_0)_\perp) e^{-\widetilde{\mathcal{V}_{L}}f(\vx_0 + s  \Theta_0 + t (\Theta_0)_\perp,\theta_0)}\dd t.
\]
Let us write $H(s,t)$ for the integrand in the previous formula. Then $H \in H^{3/2+,2}(\mathbb{R}^2)$ and $\frac{\partial H}{\partial s} \in H^{1/2+,2}(\mathbb{R}^2)$ and by Theorem \ref{trace_thm} we can restrict to the segment of integration and the result is in $H^{0+,2}$ on this segment. Thus, for $s$ larger than $s_0$, we have
\[
R'(s) = t_+'(s) H(s,t_+(s)) - t'_-(s) H(s,t_-(s)) + G(s,t)
\]
where $G$ is bounded. Because the curvature of $\partial \Omega$ is positive at $\vx_0$, $t_\pm(s) \rightarrow 0$ as $s \rightarrow s_0^+$ and
\[
\lim_{s\rightarrow s_0^+} (s-s_0)^{1/2} t'_{\pm}(s) = \pm \kappa
\]
for some $\kappa > 0$. Since $H$ is continuous and $(s-s_0)^{1/2}$ is not in $H^{1,2}(\mathbb{R})$ near $s = s_0$, therefore $R$ is not in $H^{1,2}(\mathbb{R})$ near $s_0$ which completes the proof.
\end{proof}

\end{document}
